% Chapter 7: The Classification of Finite Simple Groups
\let\LocalMacro\relax

\chapter{The Classification of Finite Simple Groups}

\section{The Problem}

\subsection{Informal}
A \emph{simple group} is a nontrivial group whose only normal subgroups are the trivial group and itself. Simple groups are the ``building blocks'' of all finite groups, analogous to prime numbers for integers.

The classification problem: \emph{List all finite simple groups.}

The answer, proved over 40 years by hundreds of mathematicians, is that every finite simple group belongs to one of the following families or is one of 26 exceptional groups:

\begin{enumerate}
\item \textbf{Cyclic groups} $C_p$ of prime order $p$
\item \textbf{Alternating groups} $A_n$ for $n \geq 5$
\item \textbf{Groups of Lie type} (various infinite families)
\item \textbf{26 sporadic groups} (exceptional groups not fitting any family)
\end{enumerate}

\begin{historicalbox}
The classification is the largest theorem in mathematics, spanning approximately 15,000 pages across 500+ journal articles. It was declared complete in 2004, though a second-generation streamlined proof is still underway.
\end{historicalbox}

\subsection{Formal}


\section{Solution}

Who solved it and what techniques were required.

\section{Mathematical Theory}


\subsection{Prerequisites}
This chapter assumes familiarity with:
\begin{itemize}
\item Group theory (normal subgroups, simple groups, quotient groups)
\item Ring theory (fields, matrix groups)
\item Basic representation theory
\end{itemize}

\subsection{The Families}

\begin{definition}[Cyclic Group of Prime Order]
For each prime $p$, the cyclic group $C_p = \mathbb{Z}/p\mathbb{Z}$ is a simple group. These are the \emph{abelian simple groups}.
\end{definition}

\begin{definition}[Alternating Group]
The \emph{alternating group} $A_n$ is the group of even permutations of $n$ elements. $A_n$ is simple for all $n \geq 5$ \cite{galois1832}.
\end{definition}

\begin{definition}[Groups of Lie Type]
The \emph{groups of Lie type} are constructed from algebraic groups over finite fields. Major families include:
\begin{itemize}
\item $PSL_n(q)$: Projective special linear groups
\item $PSU_n(q)$: Projective special unitary groups
\item $PSp_{2n}(q)$: Projective symplectic groups
\item $P\Omega_{2n+1}(q)$ and $P\Omega_{2n}^{\pm}(q)$: Orthogonal groups
\item Exceptional groups: $G_2(q)$, $^2F_4(q)$, $^3D_4(q)$, $E_6(q)$, $E_7(q)$, $E_8(q)$, $F_4(q)$
\end{itemize}
\end{definition}

The groups of Lie type are parametrized by a prime power $q = p^f$ and a Dynkin diagram type. They arise as the finite quotients of algebraic groups.

\subsection{The Sporadic Groups}

The 26 sporadic groups are \cite{conway1985}:

\begin{enumerate}
\item Mathieu groups: $M_{11}, M_{12}, M_{22}, M_{23}, M_{24}$
\item Conway groups: $Co_1, Co_2, Co_3$
\item McLaughlin group: $McL$
\item Higman--Sims group: $HS$
\item Held group: $He$
\item Rudvalis group: $Ru$
\item Suzuki group: $Su$
\item Janko groups: $J_1, J_2, J_3, J_4$
\item Fischer groups: $Fi_{22}, Fi_{23}, Fi_{24}'$
\item Baby Monster: $B$
\item Monster (Fischer--Griess): $\mathbb{M}$
\end{enumerate}

The largest sporadic group is the Monster $\mathbb{M}$, with order approximately $8 \times 10^{53}$. Its discovery by Fischer and Griess (1982) led to the surprising \emph{Monstrous Moonshine} phenomenon linking the Monster to modular forms \cite{conway1979}.

\section{The Discovery}

\subsection{Completion of the Classification (2004)}

The classification theorem was announced as complete in 2004, when the remaining cases (primarily those involving groups of characteristic 2) were settled by Michael Aschbacher and Stephen Smith \cite{aschbacher2004}:

\begin{theorem}[Aschbacher \& Smith, 2004 \cite{aschbacher2004}]
The classification of finite simple groups is complete. Every finite simple group is one of the following:
\begin{enumerate}
\item A cyclic group of prime order
\item An alternating group $A_n$ ($n \geq 5$)
\item A group of Lie type
\item One of the 26 sporadic groups
\end{enumerate}
\end{theorem}

The final case was the classification of \emph{quasi-thin groups} (groups with certain properties related to the 2-local structure). Aschbacher and Smith's proof spanned over 2000 pages.

\subsection{The Second Generation Project}

After the classification was declared complete, a project began to rewrite the proofs in a more coherent and accessible form:

\begin{itemize}
\item \textbf{Goal}: Produce a streamlined proof of approximately 5,000 pages (down from 15,000).
\item \textbf{Leaders}: Aschbacher, Smith, Gorenstein (deceased), Lyons, Solomon.
\item \textbf{Status (2025)}: Significant progress has been made, but the project is ongoing. The revised proof is expected to be published in a series of volumes by the AMS.
\end{itemize}

\section{Impact}

\begin{itemize}
\item \textbf{Unification}: The classification unified dozens of disparate results in group theory into a single coherent picture.
\item \textbf{Sporadic groups}: The discovery of the 26 sporadic groups (especially the Monster) led to unexpected connections with other areas (moonshine, string theory).
\item \textbf{Applications}: The classification has applications to coding theory, cryptography, and the symmetry of mathematical objects.
\item \textbf{Chemistry}: The classification helps explain molecular symmetry in chemistry.
\end{itemize}

\begin{highlightbox}
The classification of finite simple groups is often called the ``Big Theorem'' of group theory. It is the largest theorem ever proved, involving the work of approximately 100 mathematicians over 40 years.
\end{highlightbox}

\section{Exercises}

\begin{exercise}[Basic]
Prove that $A_5$ is simple. \emph{Hint: Show that any nontrivial normal subgroup must contain a 3-cycle, and then use conjugation to show it contains all 3-cycles.}
\end{exercise}

\begin{exercise}[Intermediate]
Compute the order of $PSL_2(7)$. Show it equals 168 and is simple. \emph{Hint: Use the fact that $PSL_n(q)$ is simple for $n \geq 2$ except $(n,q) = (2,2)$ and $(2,3)$.}
\end{exercise}

\begin{exercise}[Challenge]
The Monster group has order approximately $8 \times 10^{53}$. How many digits does this number have? Express your answer in terms of logarithms.
\end{exercise}

\begin{exercise}
Explain why the classification of finite simple groups is analogous to the periodic table of elements in chemistry. What are the ``elements'' in each case?
\end{exercise}


\begin{exercise}[Computational]
Write a Python/SageMath script to explore a concept from this chapter. See the appendix for starter code.
\end{exercise}

\section{Timeline}

\begin{tabularx}{\linewidth}{lX}
\toprule
\textbf{Year} & \textbf{Event} \\ \midrule
1832 & Galois proves $A_n$ is simple for $n \geq 5$ \cite{galois1832} \\ 
1872 & Klein discovers $PSL_2(7)$ (simple group of order 168) \cite{klein1872} \\ 
1965 & Suzuki classifies simple groups with a specific involution structure \cite{suzuki1965} \\ 
1968 & Walter classifies simple groups with abelian Sylow 2-subgroups \cite{walter1968} \\ 
1972 & Thompson begins the classification program with N-group theory \cite{thompson1972} \\ 
1980s & Major advances: Aschbacher, Gorenstein, Lyons, Solomon \cite{gls1994} \\ 
1981 & Harada identifies the Harada--Norton group as sporadic \cite{harada1981} \\ 
1998 & Solomon completes the character table of the Baby Monster \cite{solomon1998} \\ 
2004 & Aschbacher and Smith complete the quasi-thin case \cite{aschbacher2004} \\ 
2004--2025 & Second generation project streamlines the proofs \cite{gls1994} \\ \bottomrule
\end{tabularx}

\section{Citations}

\begin{enumerate}
\item Aschbacher, M., \& Smith, S. (2004). ``The classification of quasi-thin groups.'' In Handbook of Finite Simple Groups.
\item Gorenstein, D., Lyons, R., \& Solomon, R. (1994--1998). ``The classification of the finite simple groups.'' Mathematical Surveys and Monographs, Volumes 1--3, AMS.
\item Wilson, R. A. (2009). ``The finite simple groups.'' Graduate Texts in Mathematics, 251, Springer.
\item Aschbacher, M. (2004). ``The classification of finite simple groups.'' In Handbook of Finite Simple Groups.
\item Conway, J. H., \& Norton, S. P. (1979). ``Monstrous moonshine.'' Bulletin of the London Mathematical Society, 11(3), 308--339.
\item Conway, J. H., et al. (1985). ``Atlas of Finite Groups.'' Oxford University Press.
\item Galois, E. (1832). ``Letter to Auguste Chevalier.'' (Published posthumously).
\item Klein, F. (1872). ``On the order-seven transformation of elliptic functions.'' Mathematische Annalen.
\item Suzuki, M. (1965). ``On finite groups with a partition.'' Journal of Mathematics of Osaka City University.
\item Walter, J. H. (1968). ``The characterization of finite groups with abelian Sylow 2-subgroups.'' Annals of Mathematics.
\item Thompson, J. G. (1972). ``N-groups and the classification of finite simple groups.'' Bulletin of the AMS.
\item Harada, K. (1981). ``On the simple group $F$ of order $2^{14} \cdot 3^6 \cdot 5^6 \cdot 7 \cdot 11 \cdot 19$.'' Proceedings of the Japan Academy.
\item Solomon, R. (1998). ``The character table of the Baby Monster.'' Journal of Algebra.
\end{enumerate}